\specialchapter{SUMMARY OF THE PRINCIPAL PROPERTIES OF ROOT SYSTEMS}

(In this summary we restrict ourselves to the case of reduced root systems and work over the field of real numbers.)

\begin{enumerate}
\def\labelenumi{\arabic{enumi})}
\tightlist
\item
  Let V be a real vector space. A reduced root system in V is a subset R of V with the following properties:
\end{enumerate}

\begin{enumerate}
\def\labelenumi{(\roman{enumi})}
\item
  R is finite and generates V.
\item
  For all \(\alpha \in \mathbb{R}\) , there exists \(\alpha^{\vee} \in V^{*}\) such that \(\langle \alpha, \alpha^{\vee} \rangle = 2\) . and such that the map
\end{enumerate}

\[
s _ {\alpha}: x \mapsto x - \langle x, \alpha^{\vee} \rangle \alpha
\]

from V to V transforms R to R.

\begin{enumerate}
\def\labelenumi{(\roman{enumi})}
\setcounter{enumi}{2}
\item
  For all \(\alpha \in \mathbb{R}\) . \(\alpha^{\vee}(\mathbb{R}) \subseteq \mathbf{Z}\) .
\item
  If \(\alpha \in \mathbb{R}\) , then \(2\alpha \notin \mathbb{R}\) .
\end{enumerate}

In view of (i), the element \(\alpha^{\vee}\) , whose existence is guaranteed by (ii), is unique; thus (iii) makes sense. The map \(s_{\alpha}\) is a reflection that leaves fixed the points of \(\mathrm{L}_{\alpha} = \mathrm{Ker}(\alpha^{\vee})\) and transforms \(\alpha\) to \(-\alpha\) .

The elements of R are called roots. The dimension of V is called the rank of the root system.

\begin{enumerate}
\def\labelenumi{\arabic{enumi})}
\setcounter{enumi}{1}
\item
  The group of automorphisms of V that leave R stable is denoted by A(R). The \(s_{\alpha}\) ( \(\alpha \in R\) ) generate a subgroup W(R) of A(R), called the Weyl group of R; this subgroup is normal in A(R). The only reflections belonging to W(R) are the \(s_{\alpha}\) .
\item
  The set \(R^{-}\) of \(\alpha^{-}\) (for \(\alpha \in R\) ) is a reduced root system in \(V^{*}\) , called the inverse system of R. The map \(\alpha \mapsto \alpha^{-}\) is a bijection, called canonical, from R to \(R^{-}\) . We have \((\mathrm{R}^{-})^{-} = \mathrm{R}\) , and the canonical bijections \(R \to R^{-}\) . \(R^{-} \to R\) are mutual inverses. The map \(u \mapsto {}^{t}u^{-1}\) defines an isomorphism from \(W(R)\) to \(W(R^{-})\) by means of which we identify these two groups.
\item
  Let V be a real vector space that is the direct sum of vector subspaces \(V_{1},\ldots,V_{r}\) . For all i, let \(R_{i}\) be a reduced root system in \(V_{i}\) . Then the union R of the \(R_{i}\) is a root system in V, called the direct sum of the \(R_{i}\) . The group \(W(R)\) can be identified with the product of the \(W(R_{i})\) . If \(R \neq \varnothing\) and R is not the direct sum of two non-empty root systems. R is said to be irreducible. This is equivalent to saying that \(W(R)\) is irreducible. Every reduced root system
\end{enumerate}

R is a direct sum of irreducible reduced root systems. uniquely determined up to a permutation, and called the irreducible components of R.

\begin{enumerate}
\def\labelenumi{\arabic{enumi})}
\setcounter{enumi}{4}
\item
  Let R be a reduced root system in V. There exist scalar products on V invariant under \(W(R)\) . In the following, we denote by \((x|y)\) such a scalar product. If V and \(V^{*}\) are identified using \((x|y)\) , we have \(\alpha^{\vee} = \frac{2\alpha}{(\alpha|\alpha)}\) . The reflection \(s_{\alpha}\) is the orthogonal reflection which transforms \(\alpha\) to \(-\alpha\) . If R is irreducible, the Weyl group acts transitively on the set of roots of a given length. If R is irreducible, the scalar product \((x|y)\) is unique up to multiplication by a constant.
\item
  Let R be a reduced root system. For \(\alpha, \beta \in \mathbb{R}\) , put
\end{enumerate}

\[
\langle \alpha , \beta^{\vee} \rangle = n (\alpha , \beta) \in \mathbf {Z}.
\]

We have

\[
\begin{array}{c} n (\alpha , \alpha) = 2, \\ s _ {\beta} (\alpha) = \alpha - n (\alpha , \beta) \beta , \\ n (\alpha , \beta) = \frac {2 (\alpha | \beta)}{(\beta | \beta)}. \end{array}
\]

The only possibilities are the following, up to interchanging \(\alpha\) and \(\beta\) :

\[
n (\alpha , \beta) = n (\beta , \alpha) = 0; \quad (\alpha , \beta) = \frac {\pi}{2}; \quad s _ {\alpha} s _ {\beta} \text {   of   order   } 2;
\]

\[
n (\alpha , \beta) = n (\beta , \alpha) = 1; \quad (\alpha , \beta) = \frac {\pi}{3}; \quad s _ {\alpha} s _ {\beta} \text {   of   order   } 3;
\]

\[
n (\alpha , \beta) = n (\beta , \alpha) = - 1; \quad (\alpha , \beta) = \frac {2 \pi}{3}; \quad s _ {\alpha} s _ {\beta} \text {   of   order   } 3;
\]

\[
n (\alpha , \beta) = 1, n (\beta , \alpha) = 2; \quad (\alpha , \beta) = \frac {\pi}{4}; \quad s _ {\alpha} s _ {\beta} \text {   of   order   } 4;
\]

\[
n (\alpha , \beta) = - 1, n (\beta , \alpha) = - 2; (\alpha , \beta) = \frac {3 \pi}{4}; s _ {\alpha} s _ {\beta} \text {   of   order   } 4;
\]

\[
n (\alpha , \beta) = 1, n (\beta , \alpha) = 3; \quad (\alpha , \beta) = \frac {\pi}{6}; \quad s _ {\alpha} s _ {\beta} \text {   of   order   } 6;
\]

\[
n (\alpha , \beta) = - 1, n (\beta , \alpha) = - 3; (\widehat {\alpha , \beta}) = \frac {5 \pi}{6}; s _ {\alpha} s _ {\beta} \text {   of   order   } 6;
\]

\[
n (\alpha , \beta) = n (\beta , \alpha) = 2; \quad \alpha = \beta ; \quad n (\alpha , \beta) = n (\beta , \alpha) = - 2; \quad \alpha = - \beta ;
\]
\[
\begin{gathered}
 n (\alpha , \beta) = n (\beta , \alpha) = 1 \Longrightarrow \| \alpha \| = \| \beta \| ; \quad n (\alpha , \beta) = n (\beta , \alpha) = -1 \Longrightarrow \| \alpha \| = \| \beta \| ; \\
 (n (\alpha , \beta) , n (\beta , \alpha)) = (1,2) \Longrightarrow \| \beta \| = \sqrt {2} \| \alpha \| ; \quad (n (\alpha , \beta) , n (\beta , \alpha)) = (-1,-2) \Longrightarrow \| \beta \| = \sqrt {2} \| \alpha \| ; \\
 (n (\alpha , \beta) , n (\beta , \alpha)) = (1,3) \Longrightarrow \| \beta \| = \sqrt {3} \| \alpha \| ; \quad (n (\alpha , \beta) , n (\beta , \alpha)) = (-1,-3) \Longrightarrow \| \beta \| = \sqrt {3} \| \alpha \| .
\end{gathered}
\]

\begin{enumerate}
\def\labelenumi{\arabic{enumi})}
\setcounter{enumi}{6}
\item
  Let \(\alpha, \beta \in \mathbb{R}\) . If \((\alpha|\beta) > 0\) , \(\alpha + \beta\) is a root unless \(\alpha = \beta\) . If \((\alpha|\beta) < 0\) , \(\alpha + \beta\) is a root unless \(\alpha = -\beta\) .
\item
  Let \(\alpha, \beta\) be two non-proportional roots. The set I of \(j \in \mathbf{Z}\) such that \(\beta + j\alpha \in \mathbb{R}\) is an interval \([-q, p]\) of \(\mathbf{Z}\) containing 0. We have
\end{enumerate}

\[
p - q = - n (\beta , \alpha), \quad \frac {q + 1}{p} = \frac {(\beta + \alpha | \beta + \alpha)}{(\beta | \beta)}.
\]

Let S be the set of \(\beta + j\alpha\) for \(j \in I\) . Then \(s_{\alpha}(S) = S\) and \(s_{\alpha}(\beta + p\alpha) = \beta - q\alpha\) . We call S the \(\alpha\) -chain of roots defined by \(\beta\) , \(\beta - q\alpha\) its origin, \(\beta + p\alpha\) its end, and \(p + q\) its length.

If T is an \(\alpha\) -chain with origin \(\gamma\) , the length of T is \(-n(\gamma, \alpha)\) .

\begin{enumerate}
\def\labelenumi{\arabic{enumi})}
\setcounter{enumi}{8}
\item
  Let X be the union of the \(\operatorname{Ker}\alpha^{\vee}\) ( \(\alpha\in R\) ). The connected components of V-X are called the chambers of R in V. They are open simplicial cones. The Weyl group acts simply-transitively on the set of chambers. If C is a chamber, \(\overline{C}\) is a fundamental domain for \(\mathrm{W}(\mathrm{R})\) . We have \((x|y)>0\) for \(x,y\in C\) . The bijection from V to \(V^{*}\) corresponding to \((x|y)\) defines a bijection from the set of chambers of R in V to the set of chambers of \(R^{\vee}\) in \(V^{*}\) ; we denote by \(C^{\vee}\) the chamber that is the image of C under this bijection.
\item
  Let C be a chamber of R. Let \(L_{1}, L_{2}, \ldots, L_{l}\) be the walls of C. For all i, there exists a unique root \(\alpha_{i}\) such that \(L_{i} = L_{\alpha_{i}}\) and \(\alpha_{i}\) is on the same side of \(L_{i}\) as C. The family \((\alpha_{1}, \ldots, \alpha_{l})\) is a basis of V, and C is the set of \(x \in V\) such that \(\langle \alpha_{i}, x \rangle > 0\) for all i, in other words such that \((\alpha_{i} | x) > 0\) for all i. Then \(\{\alpha_{1}, \ldots, \alpha_{l}\}\) is called the basis B(C) of R defined by C. We have \((\alpha_{i} | \alpha_{j}) \leqslant 0\) when \(i \neq j\) . The group W(R) acts simply-transitively on the set of bases. Every root is transformed by some element of W(R) to an element of B(C). We have \(\{\alpha_{1}^{\vee}, \ldots, \alpha_{l}^{\vee}\} = B(C^{\vee})\) .
\item
  Put \(s_{\alpha_i} = s_i\) , let \(S\) be the set of the \(s_i\) , and let \(m_{ij}\) be the order of \(s_i s_j\) . The pair (W(R), S) is a Coxeter system with Coxeter matrix ( \(m_{ij}\) ): in other words, W(R) is defined by the family of generators ( \(s_i\) ) \(_{1 \leqslant i \leqslant l}\) and the relations ( \(s_i s_j\) ) \(^{m_{ij}} = 1\) . Two elements \(s_i\) and \(s_j\) are conjugate in W(R) if and only if there exists a sequence of indices ( \(i_1, i_2, \ldots, i_q\) ) such that \(i_1 = i, i_q = j\) and each of the \(m_{i_l i_{l+1}}\) is equal to 3.
\item
  Let \(n_{ij} = n(\alpha_i, \alpha_j)\) . The matrix \((n_{ij})_{1 \leqslant i,j \leqslant l}\) is called the Cartan matrix of \(R\) . It is independent (up to a permutation of \(1, 2, \ldots, l\) ) of the choice of \(C\) . We have \(n_{ii} = 2\) . \(n_{ij} \in \{0, -1, -2, -3\}\) for \(i \neq j\) . If two root systems have the same Cartan matrix they are isomorphic.
\item
  Let \(G\) be the subgroup of \(A(R)\) that leaves \(B(C)\) stable. Then \(A(R)\) is the semi-direct product of \(G\) by \(W(R)\) .
\item
  The order relation on V (resp. \(V^{*}\) ) defined by C is the order relation, compatible with the vector space structure of V (resp. \(V^{*}\) ). for which the elements \(\geqslant 0\) are the linear combinations of the \(\alpha_{i}\) (resp. \(\alpha_{i}^{\vee}\) ) with coefficients \(\geqslant 0\) . These elements are called positive for C, or for B(C). These order relations are also defined by \(C^{\vee}\) . An element of V is \(\geqslant 0\) if and only if its values on \(C^{\vee}\) are \(\geqslant 0\) . The set of elements \(\geqslant 0\) for C contains \(\overline{C}\) but is in general distinct from \(\overline{C}\) . Let \(x \in V\) . Then \(x \in \overline{C}\) if and only if \(x \geqslant w(x)\) for all \(w \in W(R)\) ; and \(x \in C\) if and only if \(x > w(x)\) for all \(w \in W(R)\) distinct from 1.
\item
  Every root is either positive or negative for \(C\) . We denote by \(R_{+}(C)\) the set of roots that are positive for \(C\) , so that \(R = R_{+}(C) \cup (-R_{+}(C))\) is a partition of R. The reflection \(s_{i}\) transforms \(\alpha_{i}\) to \(-\alpha_{i}\) and permutes the elements of \(\mathrm{R}_{+}(\mathrm{C})\) distinct from \(\alpha_{i}\) among themselves.
\item
  Let B be a basis of R. Every positive (resp. negative) root for B is a linear combination of elements of B with coefficients that are integers \(\geqslant 0\) (resp. \(\leqslant 0\) ).
\item
  Let \((\beta_{1},\beta_{2},\ldots,\beta_{n})\) be a sequence of roots that are positive for C and such that \(\beta_{1}+\cdots+\beta_{n}\) is a root. There exists a permutation \(\pi\in S_{n}\) such that, for all \(i\in\{1,2,\ldots,n\}\) , \(\beta_{\pi(1)}+\beta_{\pi(2)}+\cdots+\beta_{\pi(i)}\) is a root.
\item
  Let \(\alpha \in \mathrm{R}_{+}(\mathrm{C})\) . Then \(\alpha \in \mathrm{B}(\mathrm{C})\) if and only if \(\alpha\) cannot be written as a sum of positive roots.
\item
  Let C be a chamber, \((\alpha_{1},\alpha_{2},\ldots,\alpha_{l})\) the corresponding basis. For any subset J of \(I=\{1,2,\ldots,l\}\) , let \(W_{J}\) be the subgroup of \(W(R)\) generated by the \(s_{i}\) such that \(i\in J\) . Let \(C_{J}\) be the set of linear combinations of the \(\alpha_{j}\) with \(j\in J\) and with coefficients \textgreater0, so that \(C_{J}\) is a facet of C.
\end{enumerate}

Let \(J \subset I\) . \(g \in W(R)\) . The following conditions are equivalent:

\begin{enumerate}
\def\labelenumi{\alph{enumi})}
\item
  g leaves a point of \(C_{J}\) invariant:
\item
  g leaves every point of \(C_{J}\) invariant;
\item
  \(g\) leaves every point of \(\overline{\mathbf{C}_{\mathrm{J}}}\) invariant;
\end{enumerate}

\(d) g(C_{J}) = C_{J};\)

\begin{enumerate}
\def\labelenumi{\alph{enumi})}
\setcounter{enumi}{4}
\tightlist
\item
  \(g(\overline{\mathrm{C}_{\mathrm{J}}}) = \overline{\mathrm{C}_{\mathrm{J}}}\) ;
\end{enumerate}

\(f) g \in W_{J}\) .

Let \(J, J' \subset \{1, 2, \ldots, l\}\) and \(g, g' \in W(R)\) . The following conditions are equivalent:

\begin{enumerate}
\def\labelenumi{\alph{enumi})}
\item
  \(g(\mathrm{C}_{\mathrm{J}}) = g'(\mathrm{C}_{\mathrm{J}'})\) ;
\item
  \(g(\mathrm{C}_{\mathrm{J}})\cap g'(\mathrm{C}_{\mathrm{J}'})\neq \emptyset ;\)
\item
  \(g\mathrm{W}_{\mathrm{J}} = g^{\prime}\mathrm{W}_{\mathrm{J}^{\prime}}\)
\end{enumerate}

\(d) J = J' \text{ and } g' \in gW_J.\)

Let \(J_1, J_2, \ldots, J_r \subset I\) and \(J = J_1 \cap \cdots \cap J_r\) . Then \(W_J = W_{J_1} \cap \cdots \cap W_{J_r}\) . For all \(g \in W(R)\) , there exists \(J \subset I\) such that \(\overline{C} \cap g(\overline{C}) = \overline{C}_J\) and \(g \in W_J\) .

\begin{enumerate}
\def\labelenumi{\arabic{enumi})}
\setcounter{enumi}{19}
\tightlist
\item
  Let \(\mathbf{P}\) be a subset of \(\mathbb{R}\) . Then \(\mathbf{P}\) is said to be closed if the conditions \(\alpha \in \mathbb{P}, \beta \in \mathbb{P}, \alpha + \beta \in \mathbb{R}\) imply \(\alpha + \beta \in \mathbb{P}\) ; and \(\mathbf{P}\) is said to be parabolic if \(\mathbf{P}\) is closed and \(\mathbf{P} \cup (-\mathbf{P}) = \mathbb{R}\) . The following conditions are equivalent:
\end{enumerate}

\begin{enumerate}
\def\labelenumi{\alph{enumi})}
\item
  P is parabolic;
\item
  P is closed and there exists a chamber C such that \(P \supset R_{+}(C)\) ;
\item
  there exist a chamber C and a subset \(\Sigma\) of B(C) such that P is the union of \(R_{+}(C)\) and the set Q of roots that are linear combinations of elements of \(\Sigma\) with coefficients that are integers \(\leqslant 0\) .
\end{enumerate}

Assume that these conditions are satisfied and let \(V_{1}\) be the vector subspace of V generated by \(\Sigma\) . Then

\[
P \cap (- P) = Q \cup (- Q) = V _ {1} \cap R,
\]

and \(P \cap (-P)\) is a root system in \(V_1\) with basis \(\Sigma\) .

Let \(P', C', \Sigma'\) have analogous properties. If there exists an element of \(W(R)\) transforming P to \(P'\) , there exists an element of \(W(R)\) transforming C to \(C'\) , \(\Sigma\) to \(\Sigma'\) and P to \(P'\) .

\begin{enumerate}
\def\labelenumi{\arabic{enumi})}
\setcounter{enumi}{20}
\tightlist
\item
  Let \(\mathbf{P}\) be a subset of \(\mathbb{R}\) . The following conditions are equivalent:
\end{enumerate}

\begin{enumerate}
\def\labelenumi{\alph{enumi})}
\item
  there exists a chamber C such that \(P = R_{+}(C)\) :
\item
  P is closed, and \(\{P, -P\}\) is a partition of R.
\end{enumerate}

The chamber C is then unique.

Assume that V is equipped with the structure of an ordered vector space such that every root is either positive or negative. Let \(R_{+}\) be the set of positive roots for this structure. There exists a unique chamber C such that \(\mathrm{R}_{+} = \mathrm{R}_{+}(\mathrm{C})\) .

\begin{enumerate}
\def\labelenumi{\arabic{enumi})}
\setcounter{enumi}{21}
\item
  A subset B of R is a basis of R if and only if the elements of B are linearly independent and every root is a linear combination of elements of B with coefficients that are all \(\geqslant 0\) or all \(\leqslant 0\) .
\item
  Let \(\mathbf{P}\) be a closed subset of \(\mathbf{R}\) such that \(\mathbf{P} \cap (-\mathbf{P}) = \varnothing\) . There exists a chamber \(\mathbf{C}\) such that \(\mathbf{P} \subset R_{+}(\mathbf{C})\) .
\item
  A subset P of R is called symmetric if P = -P. Let P be a subset of R, and let \(V_{1}\) (resp. \(\Gamma\) ) be the vector subspace (resp. the additive subgroup) of V generated by P. The following conditions are equivalent:
\end{enumerate}

\begin{enumerate}
\def\labelenumi{\alph{enumi})}
\item
  P is closed and symmetric;
\item
  P is closed and P is a root system in \(V_{1}\) ;
\item
  \(\Gamma \cap \mathbb{R} = \mathbb{P}\) .
\end{enumerate}

\begin{enumerate}
\def\labelenumi{\arabic{enumi})}
\setcounter{enumi}{24}
\tightlist
\item
  Assume that R is irreducible. Let C be a chamber: put
\end{enumerate}

\[
\mathrm{B} (\mathrm{C}) = \left\{\alpha_ {1}, \dots , \alpha_ {l} \right\}.
\]

There exists a highest element of R (for the order defined by C). that is, an element \(\tilde{\alpha}=n_{1}\alpha_{1}+\cdots+n_{l}\alpha_{l}\) of R such that, for any root

\[
p _ {1} \alpha_{1} + \dots + p _ {l} \alpha_{l}.
\]

\(n_1 \geqslant p_1, \ldots, n_l \geqslant p_l\) . Then \(\tilde{\alpha} \in \overline{C}\) and \(\| \tilde{\alpha} \| \geqslant \| \alpha \|\) for every root \(\alpha\) .

\begin{enumerate}
\def\labelenumi{\arabic{enumi})}
\setcounter{enumi}{25}
\tightlist
\item
  The subgroup of V generated by R is denoted by Q(R): the elements of Q(R) are called the radical weights of R. The group Q(R) is a discrete subgroup of V of rank \(l = \dim V\) . Any basis of R is a basis of Q(R).
\end{enumerate}

The subgroup of V associated to \(Q(R^{\vee})\) is denoted by \(P(R)\) ; the elements of \(P(R)\) are called the weights of R. The group \(P(R)\) is a discrete subgroup of V of rank l containing \(Q(R)\) . The groups \(P(R)/Q(R)\) , \(P(R^{\vee})/Q(R^{\vee})\) are finite and isomorphic; their order f is called the connection index of R. With the notation of 25), the order of \(W(R)\) is \(l!n_{1}n_{2}\ldots n_{l}f\) .

The group \(A(R)\) leaves \(P(R)\) , \(Q(R)\) stable, and hence acts on \(P(R)/Q(R)\) . The group \(W(R)\) acts trivially on \(P(R)/Q(R)\) , so \(A(R)/W(R)\) acts on \(P(R)/Q(R)\) .

\begin{enumerate}
\def\labelenumi{\arabic{enumi})}
\setcounter{enumi}{26}
\tightlist
\item
  Let C be a chamber. Let \(B = (\alpha_{1}, \ldots, \alpha_{l})\) be the corresponding basis of R. The dual basis \((\omega_{1}, \ldots, \omega_{l})\) of \((\alpha_{1}^{\vee}, \ldots, \alpha_{l}^{\vee})\) is a basis of P(R). The \(\omega_{i}\) are called the fundamental weights (for C, or for B). The set of linear combinations of the \(\omega_{i}\) with coefficients \textgreater{} 0 (resp. \(\geqslant 0\) ) is C (resp. \(\overline{C}\) ). The linear combinations of the \(\omega_{i}\) with integer coefficients \(\geqslant 0\) are called the dominant weights. Every element of P(R) can be transformed by some element of W(R) to a unique dominant weight. The dominant weights are the elements \(\omega\) of V such that \(\frac{2(\omega|\alpha_{i})}{(\alpha_{i}|\alpha_{i})}\) is an integer \(\geqslant 0\) for all i.
\end{enumerate}

\[
28) \text {   Let   } \rho = \frac {1}{2} \sum_ {\alpha \in R _ {+} (C)} \alpha . \text {   Then,   } \rho = \omega_ {1} + \dots + \omega_ {l} \in C.
\]

\begin{enumerate}
\def\labelenumi{\arabic{enumi})}
\setcounter{enumi}{28}
\item
  Let \(\mathbf{T}\) be the group of translations of \(\mathbf{V}^*\) whose vectors belong to \(\mathrm{Q}(\mathbb{R}^{\vee})\) . The group of affine transformations of \(\mathbf{V}^*\) generated by \(\mathbf{T}\) and \(\mathrm{W}(\mathbb{R})\) is the semi-direct product of \(\mathrm{W}(\mathbb{R})\) by \(\mathbf{T}\) . This group is called the affine Weyl group of \(\mathbb{R}\) and is denoted by \(\mathrm{W}_a(\mathbb{R})\) . It acts properly on \(\mathbf{V}^*\) . For \(\alpha \in \mathbb{R}\) and \(\lambda \in \mathbf{Z}\) , let \(s_{\alpha, \lambda}\) be the map \(x^* \mapsto x^* - \langle x^*, \alpha \rangle \alpha^{\vee} + \lambda \alpha^{\vee}\) ; this is an affine reflection, and the set \(L_{\alpha, \lambda}\) of its fixed points is defined by the equation \(\langle x^*, \alpha \rangle = \lambda\) ; we have \(L_{\alpha, \lambda} = L_{\alpha} + \frac{1}{2}\lambda \alpha^{\vee}\) . The \(s_{\alpha, \lambda}\) are the affine reflections belonging to \(W_a(\mathbb{R})\) , and they generate the group \(W_a(\mathbb{R})\) .
\item
  Let \(\mathbf{E}\) be the union of the \(L_{\alpha, \lambda}\) for \(\alpha \in \mathbb{R}\) and \(\lambda \in \mathbf{Z}\) . The connected components of \(V^{*} - E\) are called the alcoves of \(\mathbb{R}\) . If \(\mathbb{R}\) is irreducible, each alcove is an open simplex: in general, an alcove is a product of open simplices. The group \(W_{a}(\mathbb{R})\) acts simply-transitively on the set of alcoves. If \(\mathbb{C}\) is an alcove, \(\overline{\mathbb{C}}\) is a fundamental domain for \(W_{a}(\mathbb{R})\) . Let \(\sigma_{1}, \ldots, \sigma_{q}\) be the reflections in \(W_{a}(\mathbb{R})\) corresponding to the walls of \(\mathbb{C}\) ; let \(\mu_{ij}\) be the order of \(\sigma_{i}\sigma_{j}\) . Then \(W_{a}(\mathbb{R})\) is defined by the generators \(\sigma_{i}\) and the relations \((\sigma_{i}\sigma_{j})^{\mu_{ij}} = 1\) .
\item
  If \(p \in \mathrm{P}(\mathbb{R}^{\vee})\) , there exists an alcove \(C\) such that \(p\) is an extremal point of \(\overline{\mathbb{C}}\) . If \(\mathbf{C}'\) is an alcove, there is a unique radical weight that is an extremal point of \(\overline{\mathbb{C}'}\) .
\end{enumerate}

Let \(x^{*} \in V^{*}\) ; the following conditions are equivalent:

\begin{enumerate}
\def\labelenumi{\alph{enumi})}
\item
  \(x^{*} \in \mathrm{P}(\mathrm{R}^{\vee})\) ;
\item
  for all \(\alpha \in \mathbb{R}\) , the hyperplane parallel to \(\mathbf{L}_{\alpha}\) and passing through \(x^{*}\) is one of the \(\mathbf{L}_{\alpha,\lambda}\) .
\end{enumerate}

Let \(C'\) be a chamber of \(R^{\vee}\) . There exists a unique alcove \(C\) contained in \(C'\) and such that \(0 \in \overline{C}\) . Assume that \(R\) is irreducible, and let \(\beta\) be the highest root of \(R\) (for \(C'\) ): then \(C\) is the set of \(x^{*} \in C'\) such that \(\langle x^{*}, \beta \rangle < 1\) .

\begin{enumerate}
\def\labelenumi{\arabic{enumi})}
\setcounter{enumi}{31}
\item
  Let S be the symmetric algebra of V. \(S^{W}\) the subalgebra consisting of the elements invariant under \(W = W(R)\) , g the order of W. \(l = \dim V\) . There exist homogeneous, algebraically independent elements \(I_{1}, I_{2}, \ldots, I_{l}\) that generate \(S^{W}\) . The \(S^{W}\) -module S has a basis consisting of g homogeneous elements. Let a be the ideal of S generated by the homogeneous elements of \(S^{W}\) of degree \textgreater0: the representation of W on S/a, induced by the representation of W on S by passage to the quotient, is isomorphic to the regular representation of W (over R).
\item
  Let \(I_1, I_2, \ldots, I_l\) be homogeneous, algebraically independent elements generating \(S^W\) . Their degrees \(k_1, k_2, \ldots, k_l\) are uniquely determined (up to order) by \(R\) . We have \(g = k_1k_2\ldots k_l\) . The number of roots is \(2\sum_{i=1}^{l}(k_i - 1)\) .
\item
  An element A of S is said to be anti-invariant under W if \(w(A) = \det(w)\) . A for all \(w \in W\) . Let \(R = R_{1} \cup (-R_{1})\) be a partition of R, and put \(\pi = \prod_{\alpha \in R_{1}} \alpha\) . The element \(\pi\) of S is anti-invariant: the anti-invariant elements of S are the elements of the form \(\pi I\) , with \(I \in S^{W}\) .
\item
  Let E be the group algebra Z[P] of the group of weights P of R. If \(p \in P\) , we denote by \(e^{p}\) the corresponding element of E. We have \(e^{p}e^{p'} = e^{p+p'}\) and the \(e^{p}\) form a basis of E. There is an action of the group W on E such that \(w(e^{p}) = e^{w(p)}\) if \(w \in W\) and \(p \in P\) . An element \(z \in E\) is said to be anti-invariant if \(w(z) = \det(w).z\) for all \(w \in W\) . For all \(z \in E\) , put \(J(z) = \sum_{w \in W} \det(w).w(z)\) . Let C be a chamber. The elements \(J(e^{p})\) , where \(p \in P \cap C\) , form a basis of the group of anti-invariant elements of E. If \(\rho\) is half the sum of the positive roots, we have:
\end{enumerate}

\[
J \left(e ^ {\rho}\right) = e ^ {\rho} \prod_ {\alpha > 0} \left(1 - e ^ {- \alpha}\right) = \prod_ {\alpha > 0} \left(e ^ {\alpha / 2} - e ^ {- \alpha / 2}\right),
\]

the products being taken over the set of roots \textgreater{} 0.

\begin{enumerate}
\def\labelenumi{\arabic{enumi})}
\setcounter{enumi}{35}
\tightlist
\item
  With the notation of 35), put
\end{enumerate}

\[
z _ {p} = \mathrm{J} (e ^ {p + \rho}) / \mathrm{J} (e ^ {\rho}) \quad \text { for } p \in \mathrm{P}.
\]

The \(z_{p}\) , for \(p \in \mathrm{P} \cap \overline{\mathrm{C}}\) , form a basis of the group \(\mathbf{E}^{\mathbf{W}}\) of elements of E invariant under W. If \(\omega_{1}, \ldots, \omega_{l}\) are the fundamental weights of R, the elements \(z_{\omega_{i}}\) , \(1 \leqslant i \leqslant l\) , are algebraically independent and generate the ring \(\mathbf{E}^{\mathbf{W}}\) .

\begin{enumerate}
\def\labelenumi{\arabic{enumi})}
\setcounter{enumi}{36}
\tightlist
\item
  Let C be a chamber of R, \((\alpha_{1},\ldots,\alpha_{l})\) the corresponding basis. The element \(c=s_{1}s_{2}\ldots s_{l}\) of W is called the Coxeter transformation of R. The conjugacy class of c does not depend on C or on the numbering of the \(\alpha_{i}\) . The order h of c is called the Coxeter number of R. The eigenvalues of c are of the form \(\exp\frac{2i\pi m_{j}}{h}\) , where the integers \(m_{1},m_{2},\ldots,m_{l}\) (called the exponents of R) are such that \(1\leqslant m_{1}\leqslant m_{2}\leqslant\cdots\leqslant m_{l}\leqslant h-1\) .
\end{enumerate}

Assume that R is irreducible. Then

\[
m _ {1} = 1, m _ {l} = h - 1.
\]

\[
m _ {j} + m _ {l + 1 - j} = h \quad (1 \leqslant j \leqslant l).
\]

\[
m _ {1} + m _ {2} + \dots + m _ {l} = \frac {1}{2} l h = \frac {1}{2} \operatorname{Card} (\mathrm{R}).
\]

Every \(m \in \{1, 2, \ldots, h - 1\}\) coprime to h is equal to exactly one of the \(m_j\) . The numbers \(m_1 + 1, m_2 + 1, \ldots, m_l + 1\) coincide, up to order, with the integers denoted by \(k_1, k_2, \ldots, k_l\) in 33). With the notation of 25), \(n_1 + \cdots + n_l = h - 1\) . There exist l orbits of \(\{1, c, c^2, \ldots, c^{h-1}\}\) on R, each of which has h elements.

If \(h\) is even, \(c^{h/2}\) transforms \(C\) to \(-C\) . We have \(-1 \in W\) if and only if the exponents of \(W\) are all odd: in that case, \(h\) is even and \(c^{h/2} = -1\) .

\pandocbounded{\includegraphics[keepaspectratio]{images/part002_b04bf376.jpg}}
